\documentclass[12pt]{article}
\usepackage[utf8]{inputenc}
\usepackage{graphicx}
\usepackage{amsmath}
\usepackage{amssymb}
\usepackage{mathrsfs}
\usepackage{pgfornament}
\usepackage[many]{tcolorbox}
\usepackage[margin = 0.5in]{geometry}
\usepackage{collectbox}
\usepackage{mdframed}
\usepackage{xcolor}
\usepackage{mathtools}
\DeclarePairedDelimiter\ceil{\lceil}{\rceil}
\DeclarePairedDelimiter\floor{\lfloor}{\rfloor}

\newcommand\textbox[1]{%
  \parbox{.333\textwidth}{#1}%
}

\linespread{2}

\makeatletter
\newcommand{\mybox}{%
    \collectbox{%
        \setlength{\fboxsep}{2pt}%
        \fbox{\BOXCONTENT}%
    }%
}
\makeatother

\usepackage{listings}
\usepackage{color}

\definecolor{dkgreen}{rgb}{0,0.6,0}
\definecolor{gray}{rgb}{0.5,0.5,0.5}
\definecolor{mauve}{rgb}{0.58,0,0.82}

\lstset{frame=tb,
  language=R,
  aboveskip=3mm,
  belowskip=3mm,
  showstringspaces=false,
  columns=flexible,
  basicstyle={\small\ttfamily},
  numbers=none,
  numberstyle=\tiny\color{gray},
  keywordstyle=\color{blue},
  commentstyle=\color{dkgreen},
  stringstyle=\color{mauve},
  breaklines=true,
  breakatwhitespace=true,
  tabsize=3
}
\title{MAT 523 HW 4}
\begin{document}

\hfill Jan. 30th, 2019
\begin{center}
\begin{large}
\textbf{MAT 524: Functions of a Real Variable II \\
\vspace{-5mm}
Homework I \\}
\end{large}
\vspace{-2mm}
Nolan R. H. Gagnon \\
\end{center}

\noindent \textbf{{\huge [}{\Large [}} R \& F: Sec. 4.4, Pr. 28 \textbf{{\Large ]}{\huge ]}}{\Large \textbf{:}}

\vspace{3mm}

\noindent\mybox{\colorbox{yellow}{\textbf{Statement}}}\hspace{3mm}\hrulefill

\noindent Let $f$ be integrable over $E$ and let $C$ be a measurable subset of $E$. Then $$\int_{C}f = \int_{E} f\chi_C.$$

\vspace{-2mm}

\noindent\mybox{\colorbox{yellow}{\textbf{Proof}}}\hspace{3mm}\hrulefill

\noindent Since $f$ is integrable over $E$ so is $|f|$, by definition. Clearly, $|f\chi_C|\leq |f|$ on $E$, so by the Integral Comparison Test (Proposition 16), it follows that $f\chi_C$ is integrable over $E$. Thus, we can write
\begin{align}
    \int_{E} f\chi_C &= \int_{E}\left[f\chi_C\right]^+ - \int_{E}\left[f\chi_C \right]^- \\
    &= \int_{E} \max\lbrace f\chi_C, 0\rbrace - \int_{E} \max\lbrace -f\chi_C, 0\rbrace \\
   &= \int_{E} \max\lbrace f, 0\rbrace\chi_C - \int_{E} \max\lbrace -f, 0\rbrace\chi_C \\
    &= \int_{E} f^+\chi_C - \int_{E} f^-\chi_C.
\end{align}
The functions $f^+$ and $f^-$ are both non-negative, so using the theory from Section 4.3 of Royden and Fitzpatrick, we see that
\begin{align}
    \int_{E} f^+\chi_C - \int_{E} f^-\chi_C &= \int_{C} f^+ - \int_{C} f^- \\
    &= \int_{C} f.
\end{align}
Therefore, $\int_C f = \int_E f\chi_C,$ as required.

\hfill$\blacksquare$

\vfill

\begin{center}
\pgfornament[width = 50mm, color = black, symmetry = h]{83}\\
\vspace{2mm}
\vspace{-0.65cm}
\hrulefill \\
\vspace{-0.95cm}
\hrulefill
\end{center}

\pagebreak

\noindent \textbf{{\huge [}{\Large [}} R \& F: Sec. 4.4, Pr. 29 \textbf{{\Large ]}{\huge ]}}{\Large \textbf{:}}

\vspace{3mm}

\noindent\mybox{\colorbox{yellow}{\textbf{Problem}}}\hspace{3mm}\hrulefill 

\noindent For a measurable function $f$ on $[1,\infty)$ which is bounded on bounded sets, define $a_n=\int_{n}^{n+1}f$ for each natural number $n$. \textcolor[rgb]{0.1,0.5,0.4}{\textbf{(a)}} Is it true that $f$ is integrable over $[1,\infty)$ if and only if the series $\sum\limits_{n=1}^{\infty} a_n$ converges? \textcolor[rgb]{0.1,0.5,0.4}{\textbf{(b)}} Is it true that $f$ is integrable over $[1,\infty)$ if and only if the series $\sum\limits_{n=1}^{\infty} a_n$ converges absolutely?

\vspace{5mm}

\noindent\mybox{\colorbox{yellow}{\textbf{Solution}}}\hspace{3mm}\hrulefill

    \noindent \textcolor[rgb]{0.1,0.5,0.4}{\textbf{(a)}} If $f$ is integrable over $[1,\infty)$ then, by definition, $\int_{[1,\infty)} |f|$ converges. But
\begin{align}
    \int_{[1,\infty)} |f| &= \int_{1}^{2} |f| + \int_{2}^{3} |f| + \int_{3}^{4} |f| + \ldots \\
    &\geq \left|\int_{1}^{2} f \right| + \left|\int_{2}^{3} f \right| + \left|\int_{3}^{4} f \right| + \ldots \\
    &= \sum\limits_{n=1}^{\infty} |a_n|,
\end{align}
implying that $\sum\limits_{n=1}^{\infty} a_n$ converges absolutely. Thus, $\sum\limits_{n=1}^{\infty} a_n$ converges. 

As for the converse statement, it is not necessarily true. Consider the following counter-example. Let $f(x) = \pi\cos(\pi x)$. Then $a_n = \int_{n}^{n+1} \pi \cos(\pi x) \hspace{1mm} dx = \sin((n+1)\pi) - \sin(n\pi) = 0$ for all $n$. Clearly $\sum\limits_{n=1}^{\infty} a_n$ converges, but $\int_{[1,\infty)} \pi\cos(\pi x) \hspace{1mm} dx = \lim_{x\rightarrow\infty}\sin(\pi x) - \sin(\pi)$ diverges.

\vspace{5mm}

\noindent\textcolor[rgb]{0.1,0.5,0.4}{\textbf{(b)}} We already showed in the previous part that if $f$ is integrable over $[1,\infty)$ then $\sum\limits_{n=1}^{\infty} a_n$ converges absolutely. 

The converse statement also doesn't hold in this case. The same counter-example can be used here, since $\sum\limits_{n=1}^{\infty} a_n$ (with $a_n$ defined as in part (a)) also converges absolutely.

\hfill$\blacksquare$

\vfill

\begin{center}
\pgfornament[width = 50mm, color = black, symmetry = h]{83}\\
\vspace{2mm}
\vspace{-0.65cm}
\hrulefill \\
\vspace{-0.95cm}
\hrulefill
\end{center}

\end{document}
